\chapter{Future Work}\label{ch:futurework}
\graphicspath{{./figs/FutureWork/}{./figs/Aim1/}{./figs/Aim2/}{./figs/Preliminary/}}

\begin{reviewchange}
The experiments, models, datasets, and software developed in this dissertation create a set of research tools rather than a list of unfinished dissertation tasks. The BPA force model and corrected torque model provide a validated mechanical plant; the Sensory Afferent Database provides structured evidence for selecting feedback pathways; and the OpenSim-to-MuJoCo conversion, BPA interface, and SNS-Toolbox network provide a reproducible environment in which those pathways can be tested. The opportunities below describe concrete ways that researchers can reuse these contributions, compare alternatives, and extend the resulting neuromechanical system.
\end{reviewchange}

\section{A Neuromechanical Simulation Pipeline}\label{sec:fw_pipeline}

\begin{reviewchange}
Researchers can use the audited conversion and coupling workflow to compare neuromechanical controllers on a common musculoskeletal plant. A useful benchmark would freeze the OpenSim source model, converted MuJoCo model, actuator mapping, contact model, and initial posture, then publish a small suite of repeatable tasks: suspended stepping, weight-supported stepping, standing with perturbations, and treadmill walking. Each task should report joint ranges, gait period and duty factor, ground-contact timing, trunk and pelvis orientation, actuator saturation, and time to failure. Running the same saved configurations on multiple machines and software backends would separate controller behavior from numerical and implementation effects.

The conversion audit can also be generalized into a model-porting test suite. For each converted model, researchers can compare musculotendon lengths, finite-difference moment arms, passive joint behavior, maximum isometric force, and joint-torque direction against the OpenSim source over sampled postures. This would turn the dissertation's repair and audit scripts into a reusable acceptance test for future human, animal, prosthetic, or robot models rather than treating conversion as a one-time visual inspection.
\end{reviewchange}

\subsection{Alternative Mechanical Plants}

\begin{reviewchange}
The controller--plant interface can support plants other than the present MuJoCo model. Researchers could substitute a Simscape Multibody subsystem for individual joints, compare rigid-tendon and compliant-tendon actuator formulations, or replace the quasi-static BPA model with the identified pressure-dynamics model. The same neural output and sensory-channel definitions should be retained so that differences in stability, phase, and torque can be attributed to the plant formulation. The corrected static torque model is computationally inexpensive enough to serve as a feedforward term or lookup table, while the dynamic model can be used to test the effect of valve bandwidth, pressure lag, and hysteresis.
\end{reviewchange}

\subsection{Controller--Plant Coupling Experiments}

\begin{reviewchange}
The existing SNS-Toolbox coupling exposes Ia, Ib, and II proprioceptive channels, heel and toe contact signals, and per-actuator commands. Researchers can use this interface for systematic ablation studies: enable one sensory class at a time, remove a pathway after steady motion is established, reverse or scale a connection, and measure the resulting changes in phase, loading, and stability. A factorial experiment across feedback class, contact gating, perturbation type, and locomotor speed would provide stronger evidence than tuning a single successful gait. The Sensory Afferent Database can supply the biological rationale and expected sign for each tested pathway.
\end{reviewchange}

\section{A Two-Layer CPG with Sensory Feedback}\label{sec:fw_cpg}

\begin{reviewchange}
The implemented rhythm-generator and pattern-formation architecture provides a starting point for comparing spinal-control hypotheses. Researchers can hold the musculoskeletal plant and motoneuron mapping fixed while varying the organization of the rhythm generator, pattern-formation layer, commissural pathways, and inhibitory interneurons. Candidate networks should be evaluated not only by whether they produce stepping but also by rhythm robustness, basin of attraction, perturbation recovery, speed range, phase consistency, and transfer between simulation engines. Tonic-stimulation, deletion, and pathway-specific stimulation protocols from the locomotion literature provide experimentally interpretable tests before whole-body performance is optimized.
\end{reviewchange}

\subsection{Sensory Feedback and Phase Regulation}

\begin{reviewchange}
Ia, Ib, II, and contact pathways can be tested as distinct control mechanisms. Researchers can quantify how Ia velocity feedback changes swing initiation, how Ib load feedback redistributes extensor activity during stance, how II length feedback affects posture and interjoint coordination, and how heel and toe events reset or prolong the locomotor phase. The most informative experiments will compare otherwise identical networks with and without each pathway and report both successful and failed trials. These comparisons would connect the database's qualitative pathway evidence to measurable changes in the robot's kinematics and force production.
\end{reviewchange}

\subsection{Verification and Reproducibility}

\begin{reviewchange}
Each controller version should pass a staged verification sequence before gait tuning: isolated neuron and synapse tests, tonic stimulation of rhythm-generator and pattern-formation populations, left--right coordination, single-joint reflex tests, suspended whole-leg stepping, and finally weight-bearing simulation. Saved parameter files, random seeds, software versions, initial states, and quantitative pass criteria should accompany the results. This structure permits another researcher to reproduce a reported behavior and determine whether a failure arises in the neural circuit, actuator model, contact mechanics, or numerical integration.
\end{reviewchange}

\section{Extending the Sensorium: Vestibular, Ocular, and Cerebellar Contributions}\label{sec:fw_senses}

\begin{reviewchange}
The same framework can test supraspinal contributions without changing the validated actuator and joint models. Vestibular feedback can be introduced as body-orientation and angular-velocity signals and evaluated through platform rotations, pushes, and sensor bias. Visual or ocular feedback can be modeled as a slower estimate of horizon and body orientation, allowing researchers to measure whether it corrects accumulated inertial drift without destabilizing fast vestibular responses. A cerebellar-like adaptive layer can then be evaluated as a mechanism for tuning phase relationships, selecting gait-specific parameter sets, or compensating for gradual changes in actuator force.

These extensions should be assessed through controlled comparisons rather than by adding all sensing at once. For example, researchers could compare vestibular-only, visual-only, and fused estimates under identical disturbances; lesion the adaptive layer after training; or transfer learned corrections between actuator conditions. The balance-platform studies and published spinal-plus-cerebellar models provide reference behaviors, while the present simulator supplies a common plant for testing them.
\end{reviewchange}

\section{Physical Realization: A Treadmill Robot with Real-Time Neural Control}\label{sec:fw_hardware}

\begin{reviewchange}
The simulation tools define a practical route to physical experiments. A hardware implementation can divide computation between low-latency joint controllers for valves, pressure sensors, and contact sensors and a companion computer for SNS evaluation and activation commands. Researchers should first reproduce single-joint force and reflex tests, then suspended bilateral stepping, body-weight-supported treadmill stepping, and finally progressively reduced support. At each stage, the simulation and hardware should receive matched commands and report the same kinematic, pressure, force, contact, and neural variables.

The corrected force and torque models can be used online for feedforward pressure commands, safety limits, and detection of unexpected compliance or routing changes. Hardware experiments can also test alternative transmission designs, including compound-pulley arrangements that trade force for actuator travel. Comparing measured and predicted torque before and after such a change would show whether the design increases useful range of motion without moving the actuator outside its characterized operating region.
\end{reviewchange}

\section{Toward Dynamic Whole-Body Maneuvers}\label{sec:fw_kickflip}

\begin{reviewchange}
Once standing and periodic walking are reproducible, researchers can extend the benchmark in a controlled progression: gait initiation and termination, speed transitions, turning, push recovery, uneven contacts, jumping, and ballistic whole-body maneuvers. Each task requires increasingly precise coordination of contact timing, actuator saturation, body orientation, and landing preparation, making the sequence a useful stress test of the complete architecture. A kickflip remains an intentionally ambitious endpoint: not an unfinished dissertation requirement, but a compact demonstration that the actuator models, musculoskeletal design, sensory feedback, spinal pattern generation, and higher-level adaptation can be coordinated for a brief, demanding maneuver.
\end{reviewchange}
